\subsection{Whole-arrow sewing and the full labelled tower}

\begin{result}{E.42}{Proved}{Genus zero classifies the ordinary labelled completed exit system}
Let $\mathcal T$ be the system of E.16 in every stable genus,
retaining the open objects, labelled strata and boundary branches,
the actual twist-subgroup embeddings, relabeling, forgetting,
clutching, and their composition identifications.  Its continuous
equivalences are taken modulo coherent natural isomorphism.
Restriction and the extension of E.39 are mutually inverse
continuous homomorphisms
\[
 \operatorname{Aut}(\mathcal T)
       \xrightarrow{\ \operatorname{res}_0\ }\widehat{\mathrm{GT}},
 \qquad A\longmapsto[\mathcal F_A].                          \tag{E.180}
\]
In particular the genus-zero kernel fixes every crossing arrow,
up to one coherent system of object comparisons.  The proof
recovers smooth boundary sewing on the entire framed exit
categories, including their noninvertible arrows, by the
two-coordinate construction below.  It uses the full profinite
completions throughout.

The additional index-$24$ branch-order covering map of E.37
has its explicitly stated comparison with (E.180): the
parameters compatible in the outer group diagram of E.37 are exactly
$\operatorname{Fix}(\mathcal J)$.  E.44 subsequently proves
that this subgroup is all of $\widehat{\mathrm{GT}}$ through
the displayed finite-subgroup and forgetting maps.
Equality of the arithmetic image with
$\widehat{\mathrm{GT}}$ remains E.24's open equation.
\end{result}
\begin{proof}
We establish the needed functor on every Hom-space.  Let $S$
be a stable bordered surface, or a disjoint pair, with two
selected labelled parametrized boundaries $a,b$.  Write $P$
for its framed exit category.  In actual multicurve coordinates,
\[
 \Hom_P(\gamma_0,\gamma_1)=
 \{p\in\Gamma^{\partial}(S):p\gamma_0\supseteq\gamma_1\}
                                      /\Delta_{\gamma_0}.
                                                               \tag{E.181}
\]
Interior multicurves are essential, are not boundary-parallel,
and have no repeated parallel components.  For two surfaces use
the product of these categories.  Their attachment of labelled
tripods, and all the resulting framed Hom-spaces inside the
original tower, are the exact comparisons of E.40.

Attach a tripod to each of $a,b$, clutch one external leg of
each tripod, and retain their other external legs as two new,
distinct companion labels.  Call the core-interface nodes $i,j$
and the new middle node $m$, and set $T=\{i,m,j\}$.  Retain
all other boundary tags.  For each core multicurve $\gamma$,
take the actual flag
\[
 T\cup\gamma\longrightarrow\{m\}\cup\gamma
                                  \longrightarrow\gamma.
                                                               \tag{E.182}
\]
Let $L$ be the category of these flags and commutative
natural-transformation diagrams, allowing noninvertible core
generizations horizontally.  Every incidence is pinned; all
skeletal transports are those of (E.145).

The two companion-labelled tripod vertices are intrinsic in
any such diagram.  Indeed contracting a connected subgraph
$H$ to a vertex $v$ gives the exact dimension identity
\[
 3g(v)-3+\operatorname{val}(v)
  =\sum_{w\in V(H)}
       \bigl(3g(w)-3+\operatorname{val}(w)\bigr)
                        +|E(H)|.                            \tag{E.183}
\]
All summands are nonnegative by stability.  A vertex of type
$(0,3)$ therefore cannot be created by contracting an edge.
Its unique companion label identifies the same surviving
source vertex.  The edge joining the two such vertices is
$m$, and their other node edges are $i,j$.  Thus every
transporter between the specified flags fixes $i,m,j$
individually.  This also applies to self-sewing and to
nonseparating interface curves.

Restrict a diagram to the cores.  It gives a framed transporter
$p$ and one independent middle-collar displacement $n\in\mathbb Z$.
Conversely extend a framed core representative over the two
standard tripods and place $n$ positive twists on the middle
annulus.  The bottom representative modulo $\Delta_{\gamma_0}$
determines the other two coordinates by their stated twist
quotients, so these constructions are inverse on diagrams.
For the exact presentation, each tripod has three central
boundary twists.  Identify its inner twist with the associated
core twist, kill its companion twist by the ordinary marked
cap, and identify the two clutched-leg twists with $\tau_m$.
Eliminating these generators gives
\[
 L\simeq P\times B\mathbb Z_m,\qquad
 \Hom_L(\gamma_0,\gamma_1)
       =\Hom_P(\gamma_0,\gamma_1)\times\mathbb Z_m.           \tag{E.184}
\]
Composition is $(p_1,n_1)(p_0,n_0)=(p_1p_0,n_1+n_0)$.
Interior multitwists lie wholly in the first factor.  Thus
(E.184) specifies the morphisms and their composition, as
well as the groups at the objects.

Both projections in this product are retained functors.
For $q:L\to P$, omit the last arrow in (E.182).  The remaining
arrow is the actual clutching of the two input tripod flags;
normalization at its retained middle node identifies those
input factors.  It kills the middle twist.  For
$r:L\to B\mathbb Z_m$, use the middle-first ordering of the
same individually distinguished nodes, retain the interfaces
$i,j$, and normalize there.  Select the component comprising
the two middle tripods.  It is the labelled four-point sphere
with its one-node crossing.  The core generizations disappear
under this component projection, and a middle twist is the
positive meridian of that crossing.  Hence
\[
              q(p,n)=p,\qquad r(p,n)=n.                     \tag{E.185}
\]
Changing the order of the distinguished smoothing nodes uses
the actual matching flags and diamonds of E.38.  The component
selection uses the pinned clutching product and the admissible
normalization square of A.15, exactly as in E.35.  Therefore
these are operations of the retained system on the full
diagram categories.  They do not choose a normal-coordinate
section.

We verify full completion of (E.184) in the actual diagram
category.  Use the finite-pinning array comparison of E.40
to retain any core multicurve that survives in the target.
The intrinsic-node proof (E.183) is the required replacement
for the distinct-pair-label proof in E.40.  For the coarse
vertical crossing, E.29 gives the full completed diagonal
stabilizer.  For its complete adjacent flag, E.32 gives the
full completion of the same discrete stabilizer: all three
nodes are fixed individually.  The projection is the
completion of that discrete identity.  This proves the
coarsening step even when projection of a cutting group to
one open endpoint has a kernel.  For a whole block of core
deletions, the coarse fibre-product calculation (E.87), as
applied in E.40, retains the full residual twist lattice.

For clarity, on one incidence orbit write $P_i$ for the
framed endpoint groups and $S_P$ for its incidence stabilizer.
Cutting the middle chain gives endpoint groups
$P_i\times\mathbb Z_m$ and incidence group
$S_P\times\mathbb Z_m$, with the same middle coordinate in
its two endpoint maps.  E.40's compact-orbit argument therefore
gives precisely
\[
 \begin{split}
 &\frac{(\widehat P_1\times\widehat{\mathbb Z})
                \times(\widehat P_0\times\widehat{\mathbb Z})}
             {\widehat S_P\times
                      \operatorname{diag}(\widehat{\mathbb Z})}
 \ \simeq\ 
 \bigl((\widehat P_1\times\widehat P_0)/\widehat S_P\bigr)
                       \times\widehat{\mathbb Z},\\
 &[(p_1,n_1),(p_0,n_0)]\longmapsto([p_1,p_0],n_1-n_0).
 \end{split}                                                 \tag{E.186}
\]
Its inverse sets $n_0=0$, $n_1=n$.  The map is exactly $(q,r)$,
with both endpoint actions retained.  Finite incidence choices
and branch permutations are handled by the finite-pinning
comparison, so this accounts for every relevant Hom-space.
Multiplication agrees on the dense discrete products and
hence on the full compact products.

The same conclusion can be checked directly on finite
congruences.  For a finite target functor on $P\times B\mathbb Z$,
take the least common multiple of the orders of the images of
the central $\mathbb Z$ at its finitely many objects.  Its
kernel contains the product of the congruence from the
restriction to $P$ with this finite-index subgroup of
$\mathbb Z$.  These product congruences are cofinal.
Thus (E.186) is the actual comparison
$\widehat L\simeq\widehat P\times B\widehat{\mathbb Z}$.

It has a section uniquely specified on every Hom-space by
\[
 s:\widehat P\longrightarrow\widehat L,
                    \qquad qs=\mathrm{id},\quad rs=1.        \tag{E.187}
\]
In actual-object coordinates this is $(p,0)$.  Intrinsically,
take the fibre over the fixed genus-zero middle crossing,
including its specified identification.  The functor $q$
on this fibre is fully faithful and essentially surjective.
Its inverse is unique up to the unique natural isomorphism
commuting with $q$ and that identification under $r$.
This formulation applies to the full isomorphism-closed
categories, so an equivalence need not fix arbitrary
skeletal representatives.  It also rules out a displacement
on a noninvertible arrow whose endpoint-group action is fixed.

Forget the two companions in the entire flag (E.182), then
stabilize.  The three-node chain becomes the sewn node $e$,
and the exact collar map is
\[
       q_e=q_iq_mq_j,
       \qquad (u_i,u_m,u_j)\longmapsto u_i+u_m+u_j.           \tag{E.188}
\]
The first arrow of (E.182) now maps to an identity at that
node, and the last is its smoothing.  On the section (E.187)
the middle displacement is zero.  Taking the bottom horizontal
arrow gives exactly smooth sewing of the original framed core
arrow: its representative is the extension of $p$, and its
interior multitwist is the same twist in the sewn surface.
Restriction and extension give this equality on every element.
The surface inclusions preserve essential distinct interior
curves; a curve made parallel through the sewn collar would
have been parallel to one of the selected boundaries, excluded
in (E.181).  Consequently the construction recovers the
actual sewing functor on the full framed exit category.

For several sewings there is one middle coordinate for each
pair.  The section is uniquely characterized by the identity
of $q$ and the vanishing of all those $r$ coordinates.
Every order gives the same section.  Stabilization multiplies
the same collar parameters, so the operation comparisons
are associative and commute for disjoint sewings, including
self-sewings.  Retained core multicurves and their whole
generization diagrams are carried throughout.  Unpinning
uses the actual finite relabeling groupoid.  The annular
unit is its single four-point crossing from E.35, so its two
boundary twists give the same meridian without introducing
two separate interface curves.

Marked caps and ordinary forgetting are already retained.
An unmarked bordered disc cap is obtained by forgetting both
external labels of the attached rational tripod and stabilizing.
The actual map C.1 removes peripheral curves and merges
coincident ones, with the stated sums on their twist lattices.
This realizes the cap functor on the recovered framed exit
category.  For a closed target one can retain an auxiliary
boundary until the last step and then use this cap.  As in
E.39, auxiliary boundaries on all relevant components and
their relabeling torsors retain every component permutation;
the union of two auxiliary systems compares their choices.
The equalities follow from the same disjoint cap operations.

We next give explicit generators for all adjacent arrows.
Fix $\sigma=\tau\cup\{c\}\to\tau$ and normalize at $\tau$,
retaining all branch pins.  Choose a pants decomposition of
the component containing $c$ which refines $c$.  The two
pants adjacent to $c$ form a four-holed sphere; in the
self-adjacent case its support is a one-holed torus.
Smoothly sew all other pants cuffs back together, then cap
and clutch the branches of $\tau$ to retain exactly those
nodes.  This actual sequence of operations takes the
support's one-node crossing to the prescribed crossing
$\sigma\to\tau$.  The identity extension carries its
pin to the prescribed pin.  The other components contribute
identity arrows, and endpoint translates give the full biset.

The one-holed-torus crossing itself has an explicit genus-zero
source.  On the four-holed sphere, label the two sides of its
node $c$ by $\{a,a_0\}$ and $\{b,b_0\}$.  Smoothly self-sew
$a$ to $b$, and cap $b_0$ by an unmarked disc.  The operation is
\[
 (0,4)\ \xrightarrow{\ \text{sew }a,b\ }\ (1,2)
             \xrightarrow{\ \text{cap }b_0\ }\ (1,1).       \tag{E.189}
\]
After cutting $c$ in the source, its two pants join along
the new smooth collar into one genus-zero component with
boundaries $a_0,b_0,c^+,c^-$.  Capping $b_0$ leaves the
pair of pants with boundaries $a_0,c^+,c^-$.  Thus $c$
is exactly the nonseparating node of the one-holed torus.
On the target the surface is precisely the one-holed torus.
Its positive local plumbing at $c$ is untouched, so the
crossing and its meridian acquire no additional winding.
This proves that the actual four-point crossing, whole-arrow
smooth sewing, ordinary clutching and capping generate
every adjacent crossing, with its retained flags and pins.

Let now $F$ be an equivalence in the genus-zero kernel.
Choose its coherent genus-zero identification with the
identity.  E.35 and E.27 identify its action on the full
bordered sewing groupoids with the identity, with one
coherent comparison at each marked-pants object.  The
comparison at an arbitrary nodal endpoint is obtained once
and for all by normalizing its actual graph, applying these
component comparisons, and performing nodal clutching with
the fixed branch pins.  Relabeling coherence descends the
result through its finite graph-permutation groupoid.
It consequently depends only on that endpoint.

E.40 makes the full framed categories intrinsic; (E.185)--
(E.188) make their whole-arrow sewing intrinsic as well.
Every tower equivalence preserves $q,r$ and their source
and target functors.  Their unique simultaneous section
therefore gives the sewing comparison for $F$.  On endpoints
these are exactly the component capping and sewing comparisons
just chosen, so the construction introduces no comparison
depending on an incident arrow.

In these coherent endpoint coordinates the four-point
crossing is fixed.  The explicit construction (E.189) and
the preceding local generation therefore fix every adjacent
pinned crossing.  A change of pants refinement is a sequence
of pants moves in the retained cut components, whose
morphisms and operation comparisons already carry the same
identity comparison.  Thus all these constructions use the
same chosen endpoint identifications.  Equivariance fixes
every endpoint translate.  E.38 then fixes every arrow and
every composition.  Undoing the coordinates gives a natural
identity of the whole system, respecting its operation
comparisons.  This proves injectivity of (E.180).

E.39 supplies its right inverse, with its actual group law,
and hence proves the asserted isomorphism.  Continuity of
restriction follows by restriction on finite categorical
quotients.  Continuity of its inverse follows from the
finite proword formulas for the quilts and the finite
generator/biset presentation in each type; the system
topology is the limit over those types and quotients.
Finally E.37 computes its stated outer group covering comparison as
$A=\mathcal J(A)$, giving precisely the subgroup stated
after (E.180).  The proof of injectivity above never uses
an ambient general multicurve-normalizer equality or
injectivity of composition into an open crossing set.
\end{proof}
